\documentclass{osucourse}
\course{5757}
%% Imported from the sheet Math 5757 shares with 5756; the sections below are common to both.

%% Not in the university's course data, so these come
%% from the published sheet this file was imported from.
\SetCourseField{title}{Modern Mathematical Methods of Relativity I, II}
\SetCourseField{credits}{3 credits}
\SetCourseField{description}{Special relativity as moving frames; tensors, exterior algebra and exterior calculus; differentiable manifolds and space time structures; parallel transport, torsion and curvatures, metric compatibility; structure equations of differential geometry.}
\SetCourseField{prereq}{Multivariable differential calculus and linear algebra (e.g. Math 2568 and/or 5101). A physics course (e.g. Physics 133 or higher). No prior knowledge of tensor calculus is assumed. However, we do assume a mature attitude towards mathematics and physics.}

\begin{document}

\CatalogDescription
\PrerequisitesAndExclusions

\begin{Purpose}
Develop from the bottom up the fundamental mathematical concepts and methods responsible for the successes in 20th century physics, mathematics, and theoretical engineering. Thus Math 5756 concretizes these developments in terms of: a) Special Relativity as the cognitive bridge to 20th century geometry b) Multilinear algebra as a source of geometrical structures, c) Linear algebra’s marriage to multi-variable calculus d) differential geometry as a three level hierarchy characterized by its  Differential structure  Parallel transport structure (a.k.a. covariant derivative)  Metric structure e) The exterior calculus f) Cartan’s two structural equations for the various flavors of differential geometry, and their application to g) The Cartan-Misner calculus
\end{Purpose}

\begin{Text}
a) \emph{Gravitation} by C. W. Misner, K. S. Thorne, and J. A. Wheeler. b) Selections from \emph{Mathematical Methods of Classical Mechanics} by V.I. Arnold. c) Selections from \emph{Lecture Notes on Elementary Topology and Geometry} by I. M. Singer. d) Selections from \emph{Spacetime Physics}, 2\textsuperscript{nd} edition, by E. Taylor and J.A. Wheeler
\end{Text}

\begin{TopicsOutline}
Math 5756 (Autumn):

A rapid course in special relativity: spacetime geometry, event horizons and accelerated frames;

\topicsub{• tensors, metric geometry vs symplectic geometry;}

\topicsub{• exterior calculus, Maxwell field equations;}

\topicsub{• manifolds, Lie derivatives, and Hamiltonian dynamics in phase space;}

\topicsub{• parallel transport, torsion, tensor calculus;}

\topicsub{• curvature and Jacobi’s equation of geodesic deviation;}

\topicsub{• Cartan’s two structural equations, metric induced properties, and Cartan-Misner curvature calculus.}

Math 5757 (Spring):

\topicsub{•     Geodesics: Hamilton-Jacobi theory, the principle of constructive interference;}

\topicsub{•     stress-energy tensor: hydrodynamics in curved spacetime and Einstein field equations;}

\topicsub{•     The conservation laws and the Bianch identities mathematized in terms of the “Boundary of a Boundary is zero (@ @ ­ = 0)” Principle.}

\topicsub{•     Solutions to the Einstein’s field equations: stars, black holes, gravitational collapse, geometry and dynamics of the universe;}

\topicsub{•     vector harmonics, tensor harmonics, acoustic and gravitational waves in violent relativistic backgrounds.}
\end{TopicsOutline}

\end{document}
